\providecommand{\vf}[1]{\textcolor{blue}{#1}}
%% Draft v3, pass 3 (2026-11-09): Murphy paragraph shape (claim first, setting, number against control, reclaim, link),
%% standard ML terms (reconstruction, source model, decoder, development set). Numbers, equations, labels and % src
%% comments unchanged. No G2 in this chapter. Legal section left as a placeholder by the author's decision.

\chapter{Translation Hub}\label{ch:hub}

A language model's hidden states retain information about the input tokens. Inversion attacks recover that input: given the hidden states of a known model, a trained decoder reconstructs the prompt \todo{cite: vec2text, embedding and hidden-state inversion}. Each such inverter is bound to one model. It learns a single architecture, tokenizer, and vocabulary, so a new target model requires a new inverter trained on that model's states.

The Platonic Representation Hypothesis suggests that the binding is not necessary \cite{huh2024platonic}. If models trained on similar data converge toward similar representations, then a learned transformation can align their hidden states in one shared space, and decoders trained on that space serve every model that maps into it. A target model would then need only a mapping into the space, not a decoder of its own.

This chapter builds that space, the translation hub, and measures how much input information it preserves. Model-specific adapters and a shared trunk map the hidden states of three public models into the hub; model-specific projectors condition each frozen model to reconstruct the source prompt. On Natural Questions queries, cross-model reconstruction reaches a mean word-F1 of \vf{0.945}, against \vf{0.120} when the decoder receives the hub representations of a different prompt. % src: results/v6_3/hpc-final/inversion.json
The hub is therefore fixed, and Chapter~\ref{ch:attack} aligns a target model to it.

\paragraph{Hypothesis H1: shared representations support cross-model inversion.} A token-wise hub trained on three open-weight models preserves enough input information for each frozen model to reconstruct the prompt from any other model's hub representations. Section~\ref{sec:hub:eval} gives the protocol and the two aggregate criteria; per-direction results report the variation that an aggregate can hide.

The chapter follows the order of construction: architecture (Section~\ref{sec:hub:design}), training objective (Section~\ref{sec:hub:training}), evaluation protocol (Section~\ref{sec:hub:eval}), results on question prompts (Section~\ref{sec:hub:results}), the planned legal-corpus evaluation (Section~\ref{sec:hub:legal}), and discussion (Section~\ref{sec:hub:discussion}).

%% ------------------------------------------------------------------
\section{Hub Design}\label{sec:hub:design}

The hub is a shared vector space with one entrance and one exit per model (Figure~\ref{fig:hub:arch}). The source models are Qwen3.5-0.8B, OLMo-2-1B, and SmolLM2-1.7B. Each contributes hidden states from the block at three quarters of its depth: block \vf{18} of \vf{24} in Qwen (width \vf{1024}), block \vf{12} of \vf{16} in OLMo (width \vf{2048}), and block \vf{18} of \vf{24} in SmolLM (width \vf{2048}). % src: configs/v6_3.yaml (d3_block, width); configs/v3.yaml (blocks, depth_fractions)
\todo{One sentence on why this depth, with its source.}

\begin{figure}[ht]
\centering
\includegraphics[width=\textwidth]{ch3/hub_architecture}
\caption{The translation hub maps each model's hidden states into one shared space, and every frozen model decodes from it. Model-specific adapters and a shared trunk produce the 512-dimensional hub representation; model-specific projectors map it into the input space of each frozen decoder. Shaded boxes are trainable modules, dashed boxes are frozen models, and the row of cells is the hub representation, a vector with no parameters of its own.}
\label{fig:hub:arch}
\end{figure}

A hidden state enters the hub through a model-specific adapter and a shared trunk. The adapter maps the hidden state $s_t$ of model $a$ to a common width of \vf{1024}; the trunk then produces the hub representation $m_t$ of \vf{512} dimensions:
\begin{equation}
m_t = W_h\, R_0\big(R_a(W_a s_t + b_a)\big) + b_h,
\qquad
R(z) = z + V_2\,\mathrm{SiLU}\big(V_1\,\mathrm{LN}(z) + c_1\big) + c_2 .
\label{eq:hub:map}
\end{equation}
Here $W_a$ is the linear layer of model $a$'s adapter, $R_a$ its residual block, $R_0$ the residual block of the shared trunk, and $W_h$ the shared linear layer from \vf{1024} to \vf{512}. Each residual block $R$ applies layer normalization, widens to \vf{2048}, and returns to \vf{1024}. % src: src/v3/bridge.py (SourceAdapter, ResidualMLP); src/v6_2/bridge.py (H3BridgeBundle); configs/v6_3.yaml (architecture)

A hub representation leaves the hub through a model-specific affine projector, which maps $m_t$ to the decoder's input width. The projected vectors form a memory prefix with one row per source token, and the frozen language model generates the reconstructed prompt after its start token; in that role the model is called a decoder. Projectors receive unnormalized hub representations, a choice Section~\ref{sec:hub:projectors} tests.

A direction pairs a source model with a decoder, e.g., Qwen $\to$ hub $\to$ OLMo. Three models define nine directions: six cross-model directions and three self-reconstruction directions. The cross-model directions carry the claim of the chapter, because in them the decoder reconstructs from states it never saw during its own training.

Only the adapters, the trunk, and the projectors are trainable; the three language models remain frozen as both sources and decoders.

The hub maps each token independently. Equation~\ref{eq:hub:map} has no positional input and no attention across tokens, so each hub vector depends on one source hidden state, although that state may already encode prompt context. The token-wise design is what Chapter~\ref{ch:attack} relies on: a target model can be aligned from paired hidden states one row at a time, and the alignment map adds no context of its own.

%% ------------------------------------------------------------------
\section{Training}\label{sec:hub:training}

Training combines prompt reconstruction with two alignment objectives. Each prompt passes through all three frozen source models, and every update includes all nine directions. The reconstruction loss preserves input information; the two auxiliary losses align the three models inside the hub.

\paragraph{Token groups.} The three tokenizers segment a prompt differently, so token positions do not correspond across models. For each pair of models, tokens are therefore aligned by character span: two tokens from different models are linked when their spans share a non-whitespace character, and each connected component of linked tokens is a token group. Each token's weight is the number of group characters its span covers. % src: src/v3/alignment.py (connected components, a_weights, b_weights); src/v6_2/data.py (_pair_groups)
For each model, the group representation is the weighted mean of the group's token vectors, normalized to unit length.

\paragraph{Objective.} The objective combines a decoding loss, a contrastive loss, and a geometry loss:
\begin{equation}
\mathcal{L} = \mathcal{L}_{\text{dec}} + 0.1\,\mathcal{L}_{\text{InfoNCE}} + 1.0\,\mathcal{L}_{\text{geom}},
\qquad
\mathcal{L}_{\text{dec}} = \tfrac{6}{7}\,\overline{\text{CE}}_{\text{cross}} + \tfrac{1}{7}\,\overline{\text{CE}}_{\text{self}} .
\label{eq:hub:loss}
\end{equation}
% src: src/v6_2/losses.py (weighted_direction_ce: (sum cross + 1/3 sum self)/7; h3_objective); configs/v6_3.yaml (loss)

\paragraph{Decoding term.} $\mathcal{L}_{\text{dec}}$ trains the hub to preserve the information needed for reconstruction. For each direction, the decoder receives the projected source representations as a memory prefix, then its start token, then the ground-truth prompt tokens $x_1 \ldots x_T$ in its own vocabulary. Under teacher forcing it predicts $x_1$ after the start token, $x_{t+1}$ after $x_t$, and the end-of-sequence token after $x_T$. The memory prefix is thus the decoder's only source for the first token, and the target reproduces the prompt rather than continuing it. The loss averages cross-entropy over the $T+1$ predictions and then over prompts; the prefix length follows the source tokenizer and may differ from $T$. % src: v6_3/plan.md section 4; src/v3/evaluate.py (teacher_forced_batch)
The bars in Equation~\ref{eq:hub:loss} average over the six cross-model and the three self-reconstruction directions. Each self direction carries one third of a cross direction's weight, so the objective trains the hub mostly to translate.

\paragraph{Contrastive term.} $\mathcal{L}_{\text{InfoNCE}}$ places the same text at the same point of the hub, whichever model produced it. For a model pair $(a,b)$, let $g^a_i$ and $g^b_i$ be the group vectors of group $i$, computed from hub representations normalized to unit length. Group $i$ under one model must identify group $i$ under the other among the groups of the batch:
\begin{equation}
\ell^{\,a\to b}_i = -\log\frac{\exp(g^a_i\cdot g^b_i/\tau)}{\sum_{j\in\mathcal{N}_i}\exp(g^a_i\cdot g^b_j/\tau)},
\qquad \tau = 0.05 ,
\label{eq:hub:infonce}
\end{equation}
where $\mathcal{N}_i$ contains group $i$ and the eligible negatives of the batch. Negatives exclude groups with the same text as group $i$ and groups from the anchor prompt's near-duplicate cluster, the other groups of the anchor prompt included. The losses $\ell^{\,a\to b}_i$ and $\ell^{\,b\to a}_i$ are averaged within each prompt, then over prompts, then over the two orders and the three model pairs. % src: src/v3/losses.py (_symmetric_group_info_nce_encoded, _negative_mask_from_codes, _mean_within_query_codes, grouped_info_nce)

\paragraph{Geometry term.} $\mathcal{L}_{\text{geom}}$ preserves the similarity structure of a batch. For a matrix $A$ of group vectors, the Gram matrix $\Gamma(A) = AA^{\top}$ holds the pairwise inner products, and $d(A,B)$ is the mean squared difference between the off-diagonal entries of $\Gamma(A)$ and $\Gamma(B)$. Let $X_a$ and $Y_a$ hold the group vectors computed from model $a$'s hidden states and from its unnormalized hub representations. Before pooling, each vector is centered by a model-specific mean over all training rows, computed once at initialization, with separate means for hidden states and hub representations; the centered vectors are pooled into groups and normalized to unit length. The loss is
\begin{equation}
\mathcal{L}_{\text{geom}} = \tfrac{1}{2}\Big(\overline{d(X_a, Y_a)} + \overline{d(Y_a, Y_b)}\Big).
\label{eq:hub:geom}
\end{equation}
The first term, averaged over the two models of each pair, preserves each model's own similarity structure; the second, averaged over the three pairs, aligns that structure across models. Both auxiliary losses are computed once per batch, not once per direction. % src: src/v4/geometry.py; src/v6_2/train.py (_geometry_sides, compute_initial_statistics)

\paragraph{Schedule and selection.} The hub is trained for \vf{2{,}500} AdamW updates of \vf{32} prompts each, drawn from \vf{40{,}000} training prompts, with one random seed. The learning rate rises to $3\times10^{-4}$ over \vf{250} warmup updates and follows a cosine decay to zero; weight decay is \vf{0.01} on the linear weights, and gradients are clipped at norm \vf{1.0}. % src: configs/v6_3.yaml (training)
Teacher-forced Acc@1 selects the checkpoint. The decoder receives the ground-truth prefix, and Acc@1 is the fraction of prompt tokens, the end-of-sequence token excluded, that it ranks first, averaged within prompts, over the \vf{5{,}000} validation prompts, and over the six cross-model directions. A checkpoint is scored every \vf{250} updates, and the earliest within \vf{0.002} of the best is retained: on Corpus A, update \vf{2{,}250}, at Acc@1 \vf{0.954}. % src: results/v6_3/hpc-final/report.json (selected_step, cross_mean); src/v3/evaluate.py (content_accuracy_at1)
The test split is excluded from training and from selection.

% \paragraph{Cost of the token-wise design.} Mapping each token independently costs selection accuracy. A variant with a two-layer transformer inside each adapter, trained on the same prompts with the same objective, reached \vf{0.987} on the selection score against \vf{0.954} for the token-wise hub. % src: results/v6_2/hpc-final/{a,b}/report.json cross_mean (variant B 0.9870, variant A 0.9536)
% The token-wise hub is retained for the property stated in Section~\ref{sec:hub:design}, and the results below are its results.

%% ------------------------------------------------------------------
\section{Evaluation}\label{sec:hub:eval}

\paragraph{Autoregressive reconstruction.} The hub is evaluated under free-running decoding. Each decoder conditions on the projected source representations and on its own generated prefix; greedy generation starts at the start token and stops at the end-of-sequence token or at a cap of the memory length plus \vf{8} tokens. % src: configs/v6_3.yaml (cap_margin 8); src/v6_3/inversion.py (CAP_RULE, batched_greedy_decode)
Two metrics compare the reconstruction with the source prompt. Word-F1 is the harmonic mean of word precision and recall over word multisets, averaged over prompts; exact match is the fraction of reconstructions equal to their reference after Unicode and whitespace normalization. % src: src/v4/clues/production_decoder.py (word_f1); src/v6_3/inversion.py (exact_match); src/v3/evaluate.py (canonicalize_text: NFKC, collapsed whitespace)
Word-F1 ignores word order, so exact match is reported as the order-sensitive metric.

\paragraph{Why free-running decoding.} Teacher-forced accuracy can overstate reconstruction, because the ground-truth prefix carries information about the next token. Free-running decoding removes that prefix, so the reconstruction must come from the memory and the decoder's own outputs. Teacher-forced Acc@1 therefore selects the checkpoint and nothing more; free-running word-F1 evaluates H1.

\paragraph{Control.} The mismatched-prompt control replaces the source memory with the hub representations of another prompt and keeps the original reference. A decoder that ignores its memory, or that generates a typical question whatever the memory holds, scores the same under matched and mismatched memories. Donors follow a fixed cyclic assignment: prompt $i$ receives the memory of prompt $i+1$, and the last prompt that of the first.
% src: src/v6_3/inversion.py (shuffled_indices: (i + 1) mod n); src/v8/projectors.py (wrong_partner: same rule)

\paragraph{Prompt sets.} Corpus A contains Natural Questions query strings, \vf{nine} words long on average (\vf{6} to \vf{19} in the validation split), split into \vf{40{,}000} training, \vf{5{,}000} validation, and \vf{5{,}000} test prompts; near-duplicate questions never cross a split, and the test split remains held out. % src: configs/v3.yaml (n_train, n_val, n_test); word counts from the references in results/v6_3/hpc-final/inversion-records.jsonl (mean 9.08, median 9)
Two overlapping sets are evaluated. The full validation split gives the per-direction scores and the selection scores. A fixed subset of \vf{500} validation prompts is the development set: Chapter~\ref{ch:attack} evaluates its target-model mappings on it, so the hub is evaluated on it as well. % src: mixed-v8-s0b.json setting.readout.source (V7.2 development subset of the validation split)

\paragraph{Criteria for H1.} One aggregate score, G1, summarizes the hub: free-running word-F1 averaged over the six cross-model directions on the validation split. A corpus fails H1 if G1 falls below \vf{0.90} or if the mismatched-prompt control on the same prompts reaches \vf{0.20}. Both thresholds come from the evaluation plan of the legal-window hubs and were applied to Corpus A retrospectively; no threshold preceded its training run. % src: v8_3/plan.md (0.90 and control below 0.20, readiness thresholds); no threshold precedes the Corpus A run
Per-direction scores accompany every aggregate, because a mean can hide one weak decoder.

%% ------------------------------------------------------------------
\section{Results on Question Prompts}\label{sec:hub:results}

\subsection{The hub reconstructs prompts across models}\label{sec:hub:translation}

Corpus A meets both criteria for H1. On its \vf{5{,}000} validation prompts, cross-model word-F1 averages \vf{0.945}, against \vf{0.120} under the mismatched-prompt control; exact match averages \vf{0.71} and is zero under the control. % src: results/v6_3/hpc-final/inversion.json cross_mean (hub 0.9445, exact 0.7116; shuffled 0.1203, exact 0.0)
The gap of \vf{0.82} between score and control shows that the decoders reconstruct from prompt-specific hub representations, not from a prior over questions. Figure~\ref{fig:hub:directions} reports all nine directions, and Table~\ref{tab:hub:reads} places each score beside its matched control.

\begin{figure}[ht]
\centering
\includegraphics[width=0.62\textwidth]{ch3/direction_recovery_nq}
\caption{The OLMo and SmolLM decoders reconstruct every source above 0.96 word-F1; the Qwen decoder reaches 0.84--0.88. Free-running word-F1, with exact match below, by source model (rows) and decoder (columns) on the \vf{5{,}000} Corpus A validation prompts; diagonal cells are self-reconstruction. The mismatched-prompt control scores \vf{0.120} to \vf{0.121} in every direction. One hub seed.}
\label{fig:hub:directions}
\end{figure}

\begin{table}[H]
\caption{Every reconstruction score on Corpus A exceeds its matched control by at least \vf{0.81}. Free-running word-F1 with decode margin \vf{8}, one hub seed; each control is evaluated on the same prompts and directions as its score.}
\label{tab:hub:reads}
\centering
\small\setlength{\tabcolsep}{4pt}
\begin{tabular}{lllcc}
\toprule
Evaluation & Prompts & Directions & Word-F1 & Control \\
\midrule
Cross-model (G1) & \vf{5{,}000} validation & six cross & \vf{0.945} & \vf{0.120} \\ % src: v6_3 inversion.json
Self-reconstruction & \vf{5{,}000} validation & three self & \vf{0.935} & \vf{0.120} \\ % src: v6_3 inversion.json, mean of the diagonal (0.8411, 0.9691, 0.9933; shuffled 0.1210, 0.1199, 0.1203)
Cross-model & \vf{500} development & six cross & \vf{0.945} & \vf{0.121} \\ % src: mixed-v8-s0b.json readout.reference (cross_macro_mean 0.9447, wrong_macro_mean 0.1214)
OLMo source & \vf{500} development & OLMo, three decoders & \vf{0.947} & \vf{0.121} \\ % src: mixed-v8-s0b.json readout.reference per_direction_mean olmo->{qwen,olmo,smollm}: hub 0.8827, 0.9716, 0.9865; wrong 0.1214, 0.1213, 0.1209
\bottomrule
\end{tabular}
\end{table}

The scores correspond to small, local errors. Table~\ref{tab:hub:examples} shows reconstructions from OLMo $\to$ Qwen, one of the two lowest cross-model directions, where the median prompt scores \vf{0.93}. % src: results/v6_3/hpc-final/inversion-records.jsonl, olmo->qwen hub arm (median 0.933)
At \vf{0.88} the decoder changes one inflected form (``sang'' to ``sung''); at \vf{0.50} it keeps the question frame and loses the place names; under mismatched memory it generates a different, fluent question. Three selected prompts show the scale of word-F1, not the frequency of each error.

\begin{table}[H]
\caption{At an intermediate score one word changes; at a low score the names are lost. Reconstructions from OLMo $\to$ Qwen on Corpus A validation prompts, selected at high, intermediate, and low word-F1; the final row is generated under mismatched memory and scored against the original prompt.}
\label{tab:hub:examples}
\centering
\begin{tabular}{p{0.36\textwidth}p{0.36\textwidth}c}
\toprule
True prompt & Decoded & Word-F1 \\
\midrule
who played uncle in only fools and horses & who played uncle in only fools and horses & \vf{1.00} \\
who has sang most songs in the world & who has sung most songs in the world & \vf{0.88} \\
is floral park in queens or long island & is flower park in king or queen parks & \vf{0.50} \\
\midrule
who played uncle in only fools and horses & who plays the girl jinx in suicide squad \emph{(mismatched prompt)} & \vf{0.25} \\
\bottomrule
\end{tabular}
\end{table}

\subsection{The Qwen decoder reconstructs least well}\label{sec:hub:qwen}

The aggregate hides one weak decoder. Both cross-model directions into Qwen fall below the \vf{0.90} that H1 asks of the mean, OLMo $\to$ Qwen at \vf{0.881} and SmolLM $\to$ Qwen at \vf{0.874}, while the four cross-model directions into OLMo and SmolLM range from \vf{0.966} to \vf{0.990} (Figure~\ref{fig:hub:directions}). % src: results/v6_3/hpc-final/inversion.json directions
The weakness follows the decoder, not the source. Qwen's own states return through the Qwen decoder at \vf{0.841}, below both cross-model directions into it, and reach \vf{0.969} and \vf{0.990} through the other two decoders. % src: same; qwen->qwen 0.8411, qwen->olmo 0.9689, qwen->smollm 0.9896
Exact match widens the gap: \vf{0.35} to \vf{0.45} into Qwen against \vf{0.78} to \vf{0.94} into the other decoders. % src: same, exact_match
H1 therefore holds on the mean and not in every direction. The scores do not identify the cause; Section~\ref{sec:hub:projectors} shows that the projector input moves it.

\subsection{The development set reproduces the validation scores}\label{sec:hub:dev}

The \vf{500} development prompts give the same picture as the full split. Mean cross-model word-F1 is \vf{0.945} against \vf{0.121} under the control, and every direction lies within \vf{0.005} of its validation score (Table~\ref{tab:hub:reads}). % src: mixed-v8-s0b.json readout.reference per_direction_mean vs v6_3 inversion.json (largest difference 0.0045, smollm->qwen)
OLMo's hub representations, decoded by all three decoders, score \vf{0.947} against \vf{0.121}: \vf{0.883} into Qwen, \vf{0.972} into OLMo, and \vf{0.987} into SmolLM. % src: mixed-v8-s0b.json readout.reference per_direction_mean olmo->{qwen,olmo,smollm}; the Chapter 4 pipeline reads the same row at 0.946 (oct_w5c-noise.json sets.iso-r0)
Chapter~\ref{ch:attack} can therefore compare its target-model scores with the hub's on the same prompts.

\subsection{Unnormalized hub representations give the higher mean}\label{sec:hub:projectors}

Normalizing the projector input lowers mean reconstruction by \vf{0.152}. In a retraining experiment with the adapters and trunk frozen, new projectors were trained on the decoding loss alone, from identical initial weights and prompt schedules, once on unnormalized and once on unit-normalized hub representations. On the development set the unnormalized input reaches mean cross-model word-F1 of \vf{0.949} and the normalized input \vf{0.797} (95\% interval of the difference \vf{0.146} to \vf{0.159} over prompts). % src: mixed-v8-s0b.json readout.{retrained_raw,normalized}.cross_macro_mean; comparison.normalized_minus_retrained_raw; src/v8/projectors.py (same initializer and schedule)
The mean hides opposite effects: the three directions into SmolLM fall from \vf{0.98}--\vf{0.99} to \vf{0.43}--\vf{0.45}, while the two cross-model directions into Qwen rise from \vf{0.89} to \vf{0.97}--\vf{0.98}. % src: same file, per_direction_mean (into SmolLM raw 0.984, 0.988, 0.991; normalized 0.429, 0.446, 0.441; into Qwen cross raw 0.894, 0.888; normalized 0.984, 0.969) and setting.loss (decoder term alone)
The unnormalized input is kept because it wins on the mean. Decoder-specific normalization is untested, the opposite responses of Qwen and SmolLM are unexplained, and the interval reflects prompt resampling under one hub seed, not variation across hubs.

%% ------------------------------------------------------------------
\section{Legal Corpora}\label{sec:hub:legal}
%% Placeholder by the author's decision (2026-11-09). The legal paragraphs, tables and figures of draft v2 return here
%% once every legal evaluation follows the protocol of Table tab:hub:reads (same prompt sets, mismatched-prompt control on each score).

The legal corpora test H1 on text of different length and vocabulary from question prompts. A separate hub is trained per corpus with the architecture and objective of Sections~\ref{sec:hub:design} and~\ref{sec:hub:training}.

\placeholder{Legal results. For each legal corpus (legal128w and the ten-word segments): the corpus description; the rows of Table~\ref{tab:hub:reads}, each with its prompt set, the directions averaged, the decode margin, exact match, the hub seed, and the mismatched-prompt control on the same prompts; the nine-direction grid; decoded examples at a high, an intermediate, and a low score; and the assessment of H1. No legal result is claimed in this draft.}

\begin{table}[H]
\caption{Planned legal-hub evaluations, in the layout of Table~\ref{tab:hub:reads}. \todo{Complete after each evaluation and its control are run on the same prompts.}}
\label{tab:hub:legal}
\centering
\begin{tabular}{lllcc}
\toprule
Corpus & Evaluation & Prompts & Word-F1 & Control \\
\midrule
legal128w & Cross-model (G1) & & \todo{} & \todo{} \\
Ten-word segments & Cross-model (G1) & & \todo{} & \todo{} \\
\bottomrule
\end{tabular}
\end{table}

%% ------------------------------------------------------------------
\section{Discussion}\label{sec:hub:discussion}

Token identity survives the passage through a shared \vf{512}-dimensional space. Three frozen models that never trained on one another's states reconstruct a prompt at mean word-F1 \vf{0.945}, against \vf{0.120} from another prompt's representations, so H1 holds on question prompts under the aggregate criteria. The result is not uniform: both cross-model directions into Qwen stay below \vf{0.90}.

The hub gives Chapter~\ref{ch:attack} a frozen space with a measured score. Aligning a target model becomes a regression onto \vf{512}-dimensional hub vectors while the hub and decoders stay frozen, and the score of a perfect alignment is known: OLMo's own hub representations, decoded by the three decoders, reach \vf{0.947} on the development set against \vf{0.121} under the control. An target-model result on the same prompts is therefore placed against that score rather than reported alone. The token-wise design adds a third property: each mapped vector carries only the information already in its source hidden state.

The design trades generality for that frozen space. One hidden state per hub vector permits direct regression from a target model but gives up the selection accuracy of a hub with attention; one hub per corpus measures each hub on the text it is applied to but does not establish a single space across domains.

Five limitations bound the claim.
\begin{itemize}
\item \textbf{One hub seed.} All scores derive from one trained hub. Variation across hub seeds is unmeasured, including the stability of Qwen's lower scores.
\item \textbf{Selection and evaluation share prompts.} The validation split determined the checkpoint and gives G1, and the development set is a subset of it. Selection ran on teacher-forced accuracy rather than free-running reconstruction, but no reported score is held-out test performance; the test split is reserved for the final evaluation. \todo{State when the test split is read.}
\item \textbf{One control.} The mismatched-prompt control shows dependence on prompt-specific representations. It does not separate token content from token order, and word-F1 ignores order; mean cross-model exact match of \vf{0.71} is the only order-sensitive score reported.
\item \textbf{Decode margin.} Generation in this chapter is capped at the memory length plus \vf{8} tokens, whereas the standard configuration of Chapter~\ref{ch:attack} permits \vf{32}. No score of this chapter has been measured at margin \vf{32}, so a Chapter~\ref{ch:attack} result cannot yet be placed against a hub score under its own decoding settings. \todo{Evaluate the OLMo row on Corpus A at margin 32 (w46 plan, step 4).} % src: w46/plan.md Setting item 6 and Steps row 4; configs/v7_cap32.yaml
\item \textbf{Three models, no pairwise baseline.} The hub contains three source models, and Chapter~\ref{ch:attack} evaluates two target models of about one billion parameters. No pairwise inversion baseline measures the hub's cost advantage at this scale. \todo{V1/V2 parity numbers.} %% the existing records compare retrieval and parameter counts, not inversion
% records (retrieval and parameter counts, not inversion): V1 hub 13,587,456 against pairwise 46,080,000 parameters (results/v1_final_result.md); V2 hub 64.8 percent fewer parameters than pairwise (documents/v2-closeout-2026-08-04.md). No hub-against-pairwise inversion score exists for this hub.
\end{itemize}
The legal-corpus experiments of Section~\ref{sec:hub:legal} will test whether these results extend beyond Natural Questions queries.
